\documentclass[10pt,a4paper]{amsart}
\usepackage{amsmath,amssymb,amsthm,graphicx}
% Retained preamble for the Stacks Project chapter snapshot in this directory.
% The chapter's own preamble.tex is a copy of the shared Stacks Project preamble; the block
% below is sliced MECHANICALLY, line for line, from it. Nothing was added, removed or
% redefined; the wrapper class, the xr-hyper/hyperref wiring and the externaldocument list
% are omitted because the assembler supplies the document class.
%
% lines taken from preamble.tex: 3, 5, 9, 11, 12, 156, 157, 158, ... (51 lines total)
%
% package coverage restored from the chapter's own preamble (2026-10-13): xy, xy-2cell

\usepackage[all]{xy}
\xyoption{2cell}
\UseAllTwocells

\newenvironment{reference}{\comment}{\endcomment}
\newenvironment{slogan}{\comment}{\endcomment}
\newenvironment{history}{\comment}{\endcomment}
\theoremstyle{plain}
\newtheorem{theorem}[subsection]{Theorem}
\newtheorem{proposition}[subsection]{Proposition}
\newtheorem{lemma}[subsection]{Lemma}
\theoremstyle{definition}
\newtheorem{definition}[subsection]{Definition}
\newtheorem{example}[subsection]{Example}
\newtheorem{exercise}[subsection]{Exercise}
\newtheorem{situation}[subsection]{Situation}
\theoremstyle{remark}
\newtheorem{remark}[subsection]{Remark}
\newtheorem{remarks}[subsection]{Remarks}
\numberwithin{equation}{subsection}
\def\lim{\mathop{\mathrm{lim}}\nolimits}
\def\colim{\mathop{\mathrm{colim}}\nolimits}
\def\Spec{\mathop{\mathrm{Spec}}}
\def\Hom{\mathop{\mathrm{Hom}}\nolimits}
\def\Ext{\mathop{\mathrm{Ext}}\nolimits}
\def\SheafHom{\mathop{\mathcal{H}\!\mathit{om}}\nolimits}
\def\SheafExt{\mathop{\mathcal{E}\!\mathit{xt}}\nolimits}
\def\Sch{\mathit{Sch}}
\def\Mor{\mathop{\mathrm{Mor}}\nolimits}
\def\Ob{\mathop{\mathrm{Ob}}\nolimits}
\def\Sh{\mathop{\mathit{Sh}}\nolimits}
\def\NL{\mathop{N\!L}\nolimits}
\def\CH{\mathop{\mathrm{CH}}\nolimits}
\def\proetale{{pro\text{-}\acute{e}tale}}
\def\etale{{\acute{e}tale}}
\def\QCoh{\mathit{QCoh}}
\def\Ker{\mathop{\mathrm{Ker}}}
\def\Im{\mathop{\mathrm{Im}}}
\def\Coker{\mathop{\mathrm{Coker}}}
\def\Coim{\mathop{\mathrm{Coim}}}
\DeclareMathSymbol{\boxtimes}{\mathbin}{AMSa}{"02}
\def\QCohstack{\mathcal{QC}\!\mathit{oh}}
\def\Cohstack{\mathcal{C}\!\mathit{oh}}
\def\Spacesstack{\mathcal{S}\!\mathit{paces}}
\def\Quotfunctor{\mathrm{Quot}}
\def\Hilbfunctor{\mathrm{Hilb}}
\def\Curvesstack{\mathcal{C}\!\mathit{urves}}
\def\Polarizedstack{\mathcal{P}\!\mathit{olarized}}
\def\Complexesstack{\mathcal{C}\!\mathit{omplexes}}
\def\Pic{\mathop{\mathrm{Pic}}\nolimits}
\def\Picardstack{\mathcal{P}\!\mathit{ic}}
\def\Picardfunctor{\mathrm{Pic}}
\def\Deformationcategory{\mathcal{D}\!\mathit{ef}}
\usepackage{cleveref}
\crefname{originalreference}{original reference}{original references}
\Crefname{originalreference}{Original reference}{Original references}
\makeatletter
\newlabel{morphisms-lemma-unramified-at-point}{{02GF}{source}{}{sourcepointer-1}{}}
\newlabel{morphisms-lemma-unramified-at-point@cref}{{[originalreference][0][]02GF}{[1][0][]source}{}{}{}}
\makeatother
\crefname{originaldefinition}{original definition}{original definitions}
\Crefname{originaldefinition}{Original definition}{Original definitions}
\makeatletter
\newlabel{definition-unramified-schemes}{{024N}{source}{}{sourcepointer-2}{}}
\newlabel{definition-unramified-schemes@cref}{{[originaldefinition][0][]024N}{[1][0][]source}{}{}{}}
\makeatother
\crefname{originalreference}{original reference}{original references}
\Crefname{originalreference}{Original reference}{Original references}
\makeatletter
\newlabel{morphisms-definition-unramified}{{02G4}{source}{}{sourcepointer-3}{}}
\newlabel{morphisms-definition-unramified@cref}{{[originalreference][0][]02G4}{[1][0][]source}{}{}{}}
\makeatother
\crefname{originallemma}{original lemma}{original lemmas}
\Crefname{originallemma}{Original lemma}{Original lemmas}
\makeatletter
\newlabel{lemma-unramified-definition}{{039J}{source}{}{sourcepointer-4}{}}
\newlabel{lemma-unramified-definition@cref}{{[originallemma][0][]039J}{[1][0][]source}{}{}{}}
\makeatother
\crefname{originalreference}{original reference}{original references}
\Crefname{originalreference}{Original reference}{Original references}
\makeatletter
\newlabel{morphisms-lemma-diagonal-unramified-morphism}{{02GE}{source}{}{sourcepointer-5}{}}
\newlabel{morphisms-lemma-diagonal-unramified-morphism@cref}{{[originalreference][0][]02GE}{[1][0][]source}{}{}{}}
\makeatother
\crefname{originalreference}{original reference}{original references}
\Crefname{originalreference}{Original reference}{Original references}
\makeatletter
\newlabel{morphisms-lemma-differentials-diagonal}{{08S2}{source}{}{sourcepointer-6}{}}
\newlabel{morphisms-lemma-differentials-diagonal@cref}{{[originalreference][0][]08S2}{[1][0][]source}{}{}{}}
\makeatother
\crefname{originalreference}{original reference}{original references}
\Crefname{originalreference}{Original reference}{Original references}
\makeatletter
\newlabel{morphisms-lemma-differentials-relative-immersion}{{01UZ}{source}{}{sourcepointer-7}{}}
\newlabel{morphisms-lemma-differentials-relative-immersion@cref}{{[originalreference][0][]01UZ}{[1][0][]source}{}{}{}}
\makeatother
\crefname{originalreference}{original reference}{original references}
\Crefname{originalreference}{Original reference}{Original references}
\makeatletter
\newlabel{morphisms-lemma-differentials-affine}{{01UT}{source}{}{sourcepointer-8}{}}
\newlabel{morphisms-lemma-differentials-affine@cref}{{[originalreference][0][]01UT}{[1][0][]source}{}{}{}}
\makeatother
\crefname{originalreference}{original reference}{original references}
\Crefname{originalreference}{Original reference}{Original references}
\makeatletter
\newlabel{properties-lemma-locally-Noetherian}{{01OW}{source}{}{sourcepointer-9}{}}
\newlabel{properties-lemma-locally-Noetherian@cref}{{[originalreference][0][]01OW}{[1][0][]source}{}{}{}}
\makeatother
\crefname{originalreference}{original reference}{original references}
\Crefname{originalreference}{Original reference}{Original references}
\makeatletter
\newlabel{morphisms-lemma-locally-finite-type-characterize}{{01T2}{source}{}{sourcepointer-10}{}}
\newlabel{morphisms-lemma-locally-finite-type-characterize@cref}{{[originalreference][0][]01T2}{[1][0][]source}{}{}{}}
\makeatother
\crefname{originalreference}{original reference}{original references}
\Crefname{originalreference}{Original reference}{Original references}
\makeatletter
\newlabel{algebra-lemma-differentials-localize}{{00RT}{source}{}{sourcepointer-11}{}}
\newlabel{algebra-lemma-differentials-localize@cref}{{[originalreference][0][]00RT}{[1][0][]source}{}{}{}}
\makeatother
\title{Unramified at a point characterized by cotangent space, local embedding into affine space, and local iso of the diagonal}
\author{The Stacks Project Authors}
\date{stacks-project:etale.tex:local-snapshot-20261012; GNU FDL 1.2 or later}
\begin{document}
\maketitle
\noindent\textit{Editorial note.} All delivered mathematical text is the author's own, copied verbatim from the pinned Stacks Project chapter file etale.tex; no statement or proof step was rewritten, reordered or removed.\par\medskip
\noindent\textit{Editorial note.} Span recomputation by label pairing (LANE\_BRIEF 10.2): in this chapter the author's proof environment follows the statement environment as a standalone block; the spec covers the complete environment pair (statement span, proof span) as recorded in the pool candidate.\par\medskip
\ifdefined\hypertarget\hypertarget{sourcepointer-1}{}\fi
\noindent\textit{Source reference.} Original reference 02GF: sources/stacks-chow/tags.txt, source lines 3226--3226. Cross-chapter standard prerequisite cited by the author inside the counted span; it is used as a pointer and its own proof is not part of this item.\par\medskip
\ifdefined\hypertarget\hypertarget{sourcepointer-2}{}\fi
\noindent\textit{Source reference.} Original definition 024N: sources/stacks-chow/tags.txt, source lines 2784--2784. Cross-chapter standard prerequisite cited by the author inside the counted span; it is used as a pointer and its own proof is not part of this item.\par\medskip
\ifdefined\hypertarget\hypertarget{sourcepointer-3}{}\fi
\noindent\textit{Source reference.} Original reference 02G4: sources/stacks-chow/tags.txt, source lines 3215--3215. Cross-chapter standard prerequisite cited by the author inside the counted span; it is used as a pointer and its own proof is not part of this item.\par\medskip
\ifdefined\hypertarget\hypertarget{sourcepointer-4}{}\fi
\noindent\textit{Source reference.} Original lemma 039J: sources/stacks-chow/tags.txt, source lines 4274--4274. Cross-chapter standard prerequisite cited by the author inside the counted span; it is used as a pointer and its own proof is not part of this item.\par\medskip
\ifdefined\hypertarget\hypertarget{sourcepointer-5}{}\fi
\noindent\textit{Source reference.} Original reference 02GE: sources/stacks-chow/tags.txt, source lines 3225--3225. Cross-chapter standard prerequisite cited by the author inside the counted span; it is used as a pointer and its own proof is not part of this item.\par\medskip
\ifdefined\hypertarget\hypertarget{sourcepointer-6}{}\fi
\noindent\textit{Source reference.} Original reference 08S2: sources/stacks-chow/tags.txt, source lines 11115--11115. Cross-chapter standard prerequisite cited by the author inside the counted span; it is used as a pointer and its own proof is not part of this item.\par\medskip
\ifdefined\hypertarget\hypertarget{sourcepointer-7}{}\fi
\noindent\textit{Source reference.} Original reference 01UZ: sources/stacks-chow/tags.txt, source lines 2428--2428. Cross-chapter standard prerequisite cited by the author inside the counted span; it is used as a pointer and its own proof is not part of this item.\par\medskip
\ifdefined\hypertarget\hypertarget{sourcepointer-8}{}\fi
\noindent\textit{Source reference.} Original reference 01UT: sources/stacks-chow/tags.txt, source lines 2422--2422. Cross-chapter standard prerequisite cited by the author inside the counted span; it is used as a pointer and its own proof is not part of this item.\par\medskip
\ifdefined\hypertarget\hypertarget{sourcepointer-9}{}\fi
\noindent\textit{Source reference.} Original reference 01OW: sources/stacks-chow/tags.txt, source lines 2209--2209. Cross-chapter standard prerequisite cited by the author inside the counted span; it is used as a pointer and its own proof is not part of this item.\par\medskip
\ifdefined\hypertarget\hypertarget{sourcepointer-10}{}\fi
\noindent\textit{Source reference.} Original reference 01T2: sources/stacks-chow/tags.txt, source lines 2359--2359. Cross-chapter standard prerequisite cited by the author inside the counted span; it is used as a pointer and its own proof is not part of this item.\par\medskip
\ifdefined\hypertarget\hypertarget{sourcepointer-11}{}\fi
\noindent\textit{Source reference.} Original reference 00RT: sources/stacks-chow/tags.txt, source lines 1015--1015. Cross-chapter standard prerequisite cited by the author inside the counted span; it is used as a pointer and its own proof is not part of this item.\par\medskip
\begin{theorem}
\label{theorem-unramified-equivalence}
Let $Y$ be a locally Noetherian scheme.
Let $f : X \to Y$ be a morphism of schemes which is locally of finite type.
Let $x$ be a point of $X$. The following are equivalent
\begin{enumerate}
\item $f$ is unramified at $x$,
\item the stalk $\Omega_{X/Y, x}$ of the module of relative differentials
at $x$ is trivial,
\item there exist open neighbourhoods $U$ of $x$ and $V$ of $f(x)$, and a
commutative diagram
$$
\xymatrix{
U \ar[rr]_i \ar[rd] & & \mathbf{A}^n_V \ar[ld] \\
& V
}
$$
where $i$ is a closed immersion defined by a
quasi-coherent sheaf of ideals $\mathcal{I}$ such that the differentials
$\text{d}g$ for $g \in \mathcal{I}_{i(x)}$ generate
$\Omega_{\mathbf{A}^n_V/V, i(x)}$, and
\item the diagonal $\Delta_{X/Y} : X \to X \times_Y X$
is a local isomorphism at $x$.
\end{enumerate}
\end{theorem}

\begin{proof}
The equivalence of (1) and (2) is proved in
Morphisms, Lemma \ref{morphisms-lemma-unramified-at-point}.

\medskip\noindent
If $f$ is unramified at $x$, then $f$ is unramified in an open
neighbourhood of $x$; this does not follow immediately
from Definition \ref{definition-unramified-schemes} of this chapter
but it does follow from
Morphisms, Definition \ref{morphisms-definition-unramified} which we
proved to be equivalent in
Lemma \ref{lemma-unramified-definition}.
Choose affine opens $V \subset Y$, $U \subset X$
with $f(U) \subset V$ and $x \in U$, such that $f$ is
unramified on $U$, i.e., $f|_U : U \to V$ is unramified.
By Morphisms, Lemma \ref{morphisms-lemma-diagonal-unramified-morphism}
the morphism $U \to U \times_V U$
is an open immersion. This proves that (1) implies (4).

\medskip\noindent
If $\Delta_{X/Y}$ is a local isomorphism at $x$, then
$\Omega_{X/Y, x} = 0$ by
Morphisms, Lemma \ref{morphisms-lemma-differentials-diagonal}.
Hence we see that (4) implies (2).
At this point we know that (1), (2) and (4) are all equivalent.

\medskip\noindent
Assume (3). The assumption on the diagram combined with
Morphisms, Lemma \ref{morphisms-lemma-differentials-relative-immersion}
show that $\Omega_{U/V, x} = 0$. Since $\Omega_{U/V, x} = \Omega_{X/Y, x}$
we conclude (2) holds.

\medskip\noindent
Finally, assume that (2) holds. To prove (3) we may localize on
$X$ and $Y$ and assume that $X$ and $Y$ are affine.
Say $X = \Spec(B)$ and $Y = \Spec(A)$.
The point $x \in X$ corresponds to a prime $\mathfrak q \subset B$.
Our assumption is that $\Omega_{B/A, \mathfrak q} = 0$
(see Morphisms, Lemma \ref{morphisms-lemma-differentials-affine} for the
relationship between differentials on schemes and modules
of differentials in commutative algebra).
Since $Y$ is locally Noetherian and $f$ locally of finite type
we see that $A$ is Noetherian and
$B \cong A[x_1, \ldots, x_n]/(f_1, \ldots, f_m)$, see
Properties, Lemma \ref{properties-lemma-locally-Noetherian} and
Morphisms, Lemma \ref{morphisms-lemma-locally-finite-type-characterize}.
In particular, $\Omega_{B/A}$ is a finite $B$-module. Hence we
can find a single $g \in B$, $g \not \in \mathfrak q$ such that
the principal localization $(\Omega_{B/A})_g$ is zero. Hence after
replacing $B$ by $B_g$ we see that $\Omega_{B/A} = 0$ (formation
of modules of differentials commutes with localization, see
Algebra, Lemma \ref{algebra-lemma-differentials-localize}). This means that
$\text{d}(f_j)$ generate the kernel of the canonical map
$\Omega_{A[x_1, \ldots, x_n]/A} \otimes_A B \to \Omega_{B/A}$.
Thus the surjection $A[x_1, \ldots, x_n] \to B$ of $A$-algebras gives the
commutative diagram of (3), and the theorem is proved.
\end{proof}

\end{document}
